\documentclass[11pt]{article}

\usepackage[margin=0.78in]{geometry}
\usepackage[T1]{fontenc}
\usepackage[utf8]{inputenc}
\usepackage{amsmath,amssymb,amsfonts,mathtools,bm}
\usepackage{booktabs,longtable,array,multirow}
\usepackage{xcolor}
\usepackage{hyperref}
\usepackage{enumitem}
\usepackage{listings}
\usepackage{caption}
\usepackage{fancyhdr}
\usepackage{lastpage}

\hypersetup{colorlinks=true,linkcolor=blue!55!black,urlcolor=blue!55!black,citecolor=blue!55!black}
\setlist{nosep,leftmargin=2em}
\setlength{\parskip}{0.42em}
\setlength{\parindent}{0pt}
\setlength{\headheight}{14pt}
\pagestyle{fancy}
\fancyhf{}
\lhead{SC-JUNA v3 receiver specification}
\rhead{Page \thepage\ of \pageref{LastPage}}

\definecolor{codebg}{RGB}{248,248,248}
\definecolor{codeframe}{RGB}{205,205,205}
\definecolor{kw}{RGB}{25,80,160}
\definecolor{comment}{RGB}{80,110,80}
\definecolor{warnbg}{RGB}{255,248,232}
\definecolor{bluebg}{RGB}{235,244,252}

\lstdefinelanguage{Julia}{
  morekeywords={function,struct,mutable,abstract,type,where,return,for,in,if,elseif,else,end,begin,let,while,using,import,export,const,true,false,nothing,Union,Vector,Matrix,Dict,Tuple,NamedTuple,Int,Float64,ComplexF64,Bool,BitVector,Symbol,NTuple,AbstractRNG},
  sensitive=true,
  morecomment=[l]{\#},
  morestring=[b]",
}
\lstset{
  language=Julia,
  basicstyle=\ttfamily\small,
  keywordstyle=\color{kw}\bfseries,
  commentstyle=\color{comment},
  backgroundcolor=\color{codebg},
  frame=single,
  rulecolor=\color{codeframe},
  columns=fullflexible,
  keepspaces=true,
  breaklines=true,
  showstringspaces=false,
  tabsize=2
}

\newcommand{\C}{\mathbb{C}}
\newcommand{\R}{\mathbb{R}}
\newcommand{\F}{\mathbb{F}}
\newcommand{\norm}[1]{\left\lVert #1 \right\rVert}
\newcommand{\abs}[1]{\left\lvert #1 \right\rvert}
\newcommand{\eps}{\varepsilon}
\newcommand{\code}[1]{\texttt{\detokenize{#1}}}
\newcommand{\TODO}[1]{\textcolor{red!70!black}{\textbf{TODO:} #1}}

\title{SC-JUNA v3: Julia Receiver Specification}
\author{SC-JUNA receiver for Doppler-distorted OFDM}
\date{\today}

\begin{document}
\maketitle

\begin{center}
\fcolorbox{codeframe}{warnbg}{%
\begin{minipage}{0.92\linewidth}
\textbf{Replaced. Do not build from this document.}\\[0.3em]
This specification has been superseded by \texttt{scjuna\_manual\_gradient\_plan.tex},
which is the document to implement from. This specification steps the bit
latents directly (\code{state.z .-= cfg.z_stepsize .* grad_z}). The deployed
receiver uses Adam, as the manual gradient plan requires. The LLR erasure stage
described here is not present in the deployed receiver. The shared constants
here (\code{RHO_MAX}, \code{LAMBDA_P_MAX}, ramp constants) differ from the
manual gradient plan. The two documents also normalize the loss differently,
so the two sets of numbers are not comparable. Marked on 31 July 2026.
\end{minipage}}
\end{center}

\begin{center}
\fbox{\begin{minipage}{0.93\linewidth}
\textbf{Purpose.} This document specifies the Julia receiver functions,
defaults, array shapes, edge rules, file schemas, candidate selection, manual
gradients, and tests. Version 1 assumes one aligned CP-included OFDM block per
packet. Preamble acquisition is not implemented here.
\end{minipage}}
\end{center}

\tableofcontents
\clearpage

\section{What changed from v2}

The earlier implementation specification left free constants, diagnostic
thresholds, pilot holdout, LDPC file format, interleaver order, partial-FFT
windows, known-Doppler validation, and the $z$-gradient open. This v3
specification states the corresponding defaults, functions, and tests.

\begin{longtable}{@{}p{0.31\linewidth}p{0.62\linewidth}@{}}
\toprule
Problem in v2 & v3 correction \\
\midrule
Derivations competed with implementation directives & Background is reduced to one page. Every main-section equation maps to a Julia function. \\
Undefined constants such as $c_i$, $\kappa_i$, $q_{\rm sym}$, $q_{\rm dec}$ & All defaults are numerical. The undefined gates are removed. Ramps depend only on iteration index. \\
Partial-FFT window ambiguity & Only non-overlapping rectangular partition-of-unity windows are allowed in v1. Candidate $M$ values are recomputed, not merged. \\
ICI coefficient indexing ambiguity & \code{C[k, off, m]} is defined exactly with no wraparound at band edges. \\
LDPC parser and encoder unspecified & v1 uses a systematic TOML schema for sparse coordinates with $H=[A\ I]$ and encoder $p=A u$ over $\F_2$. \\
Interleaver mentioned but not specified & A deterministic stride interleaver is defined, with explicit inverse and carrier-to-bit map. \\
SC-JUNA loop less concrete than KD baselines & SC-JUNA is given in Julia-like pseudocode with exact calls and return types. \\
Known-Doppler benchmark lacked a test & A deterministic Doppler undo test is included. \\
AD choice ambiguous & Production uses a hand-derived real gradient for $z$. Zygote is only an optional gradient checker. \\
Missing numerical primitive tests & Tests are specified for OFDM, partial FFT, channel, LDPC, RPC, known Doppler, and JUNA gradients. \\
\bottomrule
\end{longtable}

\section{Receiver execution order}

The supporting types, indexing, data order, and receiver functions are stated
before the receiver sequence. For an aligned CP-included OFDM block, the
receiver follows this order:

\begin{enumerate}
\item Remove the cyclic prefix and form the full and partial FFT observations.
\item Estimate the packet diagnostics and build the candidate list.
\item Initialize the JUNA state from the fit pilots.
\item Update $C$, $W$, and $z$, then form the LLRs and run BP.
\item Score every candidate and select the result.
\end{enumerate}

\section{Version-1 scope}

\subsection{In scope}

\begin{enumerate}
\item Julia simulation of one CP-OFDM packet through a time-varying underwater acoustic channel.
\item Full-FFT, RPC, known-Doppler, fixed-JUNA, and SC-JUNA receivers.
\item Rate-$1/2$ systematic LDPC encoding and belief propagation decoding.
\item Receiver-side adaptation only. No transmitter feedback, no waterfilling, and no adaptive rate control.
\item Deterministic tests sufficient to localize implementation errors before running long Monte Carlo sweeps.
\end{enumerate}

\subsection{Out of scope for v1}

\begin{enumerate}
\item Preamble acquisition, packet detection, and timing estimation. The test harness passes an already aligned CP-included OFDM block.
\item MIMO or hydrophone array processing.
\item Nonrectangular or overlapping partial-FFT windows.
\item Non-systematic LDPC encoding.
\item Production automatic differentiation for complex variables.
\item Real-time embedded optimization.
\end{enumerate}

\subsection{Required behavior}

\begin{longtable}{@{}p{0.25\linewidth}p{0.68\linewidth}@{}}
\toprule
Topic & Required behavior \\
\midrule
Indexing & Julia arrays are 1-based. Time/sample formulas may use 0-based indices. Conversion is stated explicitly. \\
Partial FFT & Recompute \code{Ypartial} for each candidate $M$. No branch merging in v1. \\
Windows & Rectangular, non-overlapping, partition-of-unity. Each sample belongs to exactly one branch. \\
ICI edges & No circular wraparound. Skip invalid neighbors. Store zeros in unused slots. \\
Normalization & OFDM IFFT is \code{sqrt(N)*ifft(X)}. FFT demodulation is \code{fft(y)/sqrt(N)}. \\
LDPC order & JUNA latents are in original LDPC codeword order. Carrier mapping uses interleaved order. BP receives deinterleaved LLRs. \\
Randomness & Every experiment owns an \code{AbstractRNG}. No function may call global \code{rand()} in production. \\
Candidate selection & Deterministic. Ties broken by fixed receiver priority and then by lower candidate index. \\
\bottomrule
\end{longtable}

\clearpage
\section{Constants table}
\label{sec:constants}

All constants used by code must come from this table or from the experiment TOML. If an experiment overrides a value, log the override.

\begin{longtable}{@{}p{0.24\linewidth}p{0.17\linewidth}p{0.52\linewidth}@{}}
\toprule
Constant & Default & Meaning \\
\midrule
\code{EPS} & $10^{-12}$ & Division and logarithm floor. \\
\code{DEFAULT_ACTIVE_MODE} & \code{:all_non_null} & Active subcarriers are all non-null carriers. \\
\code{PILOT_SPACING} & $8$ & Comb pilot spacing in active-carrier order. \\
\code{HOLDOUT_STRIDE} & $8$ & Every eighth pilot is validation if enough pilots exist. \\
\code{MIN_HOLDOUT_PILOTS} & $4$ & If fewer than this would be held out, disable holdout and use fit pilots for residual reporting only. \\
\code{LLR_CLIP} & $12.0$ & Clip channel LLRs to $[-12,12]$. \\
\code{BP_CLIP} & $20.0$ & Clip BP internal messages to $[-20,20]$. \\
\code{BP_MAXITER} & $50$ & BP iterations. \\
\code{REL_ERASE_THRESHOLD} & $0.10$ & Erase data carrier when relative reliability is below this. \\
\code{INTERLEAVER_STRIDE} & $33$ & Default stride for permutation $\pi(i)=1+((i-1)s\bmod n_b)$; must be coprime with $n_b$. \\
\code{CANDIDATE_M} & $[4,8,16,32]$ & Candidate partial-FFT branch counts. \\
\code{KMAX} & $4$ & Maximum JUNA ICI half-bandwidth. \\
\code{JUNA_MAX_OUTER} & $80$ & Nominal JUNA outer iterations. \\
\code{JUNA_RESCUE_OUTER} & $200$ & Rescue candidate outer iterations. \\
\code{Z_STEPS_PER_OUTER} & $2$ & Manual gradient steps for $z$ per outer iteration. \\
\code{Z_STEPSIZE} & $10^{-2}$ & Initial $z$ gradient step size. \\
\code{Z_CLIP} & $8.0$ & Clamp each $z_i$ after every $z$ step. \\
\code{LAMBDA_D} & $1.0$ & Observation consistency weight. \\
\code{ETA_PILOT} & $10.0$ & Pilot anchor weight. \\
\code{RHO_MAX} & $1.0$ & Maximum combiner-symbol tie weight. \\
\code{GAMMA_C} & $10^{-2}$ & Ridge for the branch coefficients in $C$. \\
\code{GAMMA_W} & $10^{-3}$ & Combiner-energy ridge. \\
\code{GAMMA_S} & $10^{-2}$ & Combiner smoothness ridge. \\
\code{LAMBDA_P_MAX} & $0.20$ & Maximum relaxed LDPC parity weight. \\
\code{T_RHO} & $5.0$ & Rho ramp time constant. \\
\code{T_PARITY} & $10.0$ & Parity ramp time constant. \\
\code{SCORE_ALPHA} & $0.5$ & Normalized syndrome weight. \\
\code{SCORE_BETA} & $0.5$ & Normalized model residual weight. \\
\code{SCORE_CHI} & $10^{-3}$ & Normalized $\norm{C}_F^2$ weight. \\
\code{SCORE_PSI} & $10^{-3}$ & Normalized $\norm{W}_F^2$ weight. \\
\code{CRC_ENABLED} & \code{false} & CRC tie-break disabled by default. \\
\code{CRC_POLY} & \code{0x1021} & CRC-16-CCITT polynomial if enabled. \\
\code{CRC_INIT} & \code{0xFFFF} & CRC initial value if enabled. \\
\bottomrule
\end{longtable}

\subsection{Iteration ramps}

Use deterministic iteration ramps:
\begin{align}
\rho^{(t)} &= \rho_{\max}\left(1-e^{-(t+1)/T_{\rho}}\right),\\
\lambda_p^{(t)} &= \lambda_{p,\max}\left(1-e^{-(t+1)/T_p}\right).
\end{align}
There are no undefined gates $q_{\rm sym}^{(t)}$ or $q_{\rm dec}^{(t)}$ in v1.

\clearpage
\section{Repository structure and build order}

\subsection{Repository tree}

\begin{lstlisting}[language={}]
SelfCalibratingJUNA.jl/
  Project.toml
  README.md
  src/
    SelfCalibratingJUNA.jl
    Configs.jl
    Indexing.jl
    QPSK.jl
    OFDM.jl
    Pilots.jl
    Interleaver.jl
    LDPC.jl
    BPDecoder.jl
    Channel.jl
    PartialFFT.jl
    Observations.jl
    FullFFTReceiver.jl
    RPCReceiver.jl
    KnownDopplerReceivers.jl
    JUNAState.jl
    JUNAGradient.jl
    JUNASolver.jl
    SelfCalibration.jl
    SCJUNAReceiver.jl
    Simulation.jl
    Metrics.jl
  data/
    ldpc/
    golden/
    results/
  scripts/
    01_golden_tests.jl
    02_uncoded_partial_fft_reproduction.jl
    03_coded_baselines.jl
    04_scjuna_sweep.jl
    05_known_doppler_stress.jl
  test/
    runtests.jl
    test_indexing.jl
    test_qpsk.jl
    test_ofdm.jl
    test_interleaver.jl
    test_ldpc.jl
    test_channel.jl
    test_partialfft.jl
    test_known_doppler.jl
    test_rpc.jl
    test_juna_gradient.jl
    test_scjuna_selection.jl
\end{lstlisting}

\subsection{Build order}

The implementation order must follow the dependency graph below. A module is not considered done until its tests pass.

\begin{enumerate}
\item \code{Indexing.jl}: signed bin, neighbor map, offset map.
\item \code{QPSK.jl}: modulator, hard demod, LLR demapper.
\item \code{OFDM.jl}: normalized IFFT/FFT, CP add/remove.
\item \code{Pilots.jl}: comb pilots and deterministic pilot holdout.
\item \code{Interleaver.jl}: stride permutation and inverse.
\item \code{LDPC.jl}, \code{BPDecoder.jl}: systematic sparse TOML parser, encoder, BP decoder.
\item \code{Channel.jl}: deterministic channel cases and time-varying channel.
\item \code{PartialFFT.jl}: rectangular windows and partial FFT identity test.
\item \code{Observations.jl}: full and partial FFT observations.
\item \code{FullFFTReceiver.jl}; then \code{RPCReceiver.jl}; then known-Doppler receivers.
\item \code{JUNAState.jl}, \code{JUNAGradient.jl}, \code{JUNASolver.jl}.
\item \code{SelfCalibration.jl}, \code{SCJUNAReceiver.jl}.
\item \code{Simulation.jl}, sweep scripts, and plots.
\end{enumerate}

\clearpage
\section{Indexing and array shapes}
\label{sec:indexing}

\subsection{FFT-bin convention}

Julia array index $k\in\{1,\ldots,N\}$ corresponds to signed FFT bin
\begin{equation}
\kappa(k;N)=
\begin{cases}
 k-1, & 1\le k\le \lfloor N/2\rfloor+1,\\
 k-1-N, & \lfloor N/2\rfloor+2\le k\le N.
\end{cases}
\end{equation}
The physical subcarrier frequency used only in Doppler formulas that include the carrier frequency is
\begin{equation}
f_k=f_c+\kappa(k;N)\Delta f,\qquad \Delta f=F_s/N.
\end{equation}

\begin{lstlisting}
function signed_bin(k::Int, N::Int)::Int
    @assert 1 <= k <= N
    mid = fld(N, 2) + 1
    return k <= mid ? k - 1 : k - 1 - N
end

function physical_subcarrier_frequency(k::Int, N::Int, fs::Float64, fc::Float64)::Float64
    return fc + signed_bin(k, N) * fs / N
end
\end{lstlisting}

\subsection{Required array shapes}

\begin{longtable}{@{}p{0.20\linewidth}p{0.18\linewidth}p{0.54\linewidth}@{}}
\toprule
Array/type field & Shape & Meaning \\
\midrule
\code{X} & $(N)$ & Frequency-domain OFDM grid. \\
\code{x} & $(N)$ & Useful time-domain block. \\
\code{xcp} & $(N+N_{cp})$ & CP-included transmitted block. \\
\code{ycp} & $(N+N_{cp})$ & CP-included received block. \\
\code{y} & $(N)$ & CP-removed received block. \\
\code{Yfull} & $(N)$ & Full FFT output. \\
\code{Ypartial} & $(N,M)$ & Partial FFT outputs; row $k$, branch $m$. \\
\code{W} & $(N,M)$ & Combiner weights. $\hat s_k=\mathbf w_k^H\mathbf Y_k$. \\
\code{C} & $(N,2K+1,M)$ & Branch-domain local ICI coefficients. \\
\code{z} & $(n_b)$ & Real latents in original LDPC codeword order. \\
\code{s} & $(N)$ & Relaxed OFDM grid: data from $z$, pilots fixed, nulls zero. \\
\code{llr_int} & $(n_b)$ & LLRs in interleaved codeword order. \\
\code{llr} & $(n_b)$ & LLRs in original LDPC codeword order. \\
\bottomrule
\end{longtable}

\subsection{ICI coefficient offset map}

For half-bandwidth $K$, \code{C[k, off, m]} stores the coefficient from transmit carrier
\begin{equation}
q=k+off-K-1
\end{equation}
to received carrier $k$ on partial-FFT branch $m$. Invalid $q$ values are skipped. No wraparound.

\begin{lstlisting}
function neighbor_indices(k::Int, N::Int, K::Int)::Vector{Int}
    @assert 1 <= k <= N
    return collect(max(1, k - K):min(N, k + K))
end

function c_offset(k::Int, q::Int, K::Int)::Int
    off = q - k + K + 1
    @assert 1 <= off <= 2K + 1
    return off
end

function q_from_offset(k::Int, off::Int, K::Int)::Int
    @assert 1 <= off <= 2K + 1
    return k + off - K - 1
end
\end{lstlisting}

\clearpage
\section{Core Julia types}

\begin{lstlisting}
struct OFDMConfig
    N::Int
    Ncp::Int
    fs::Float64
    fc::Float64
    active_indices::Vector{Int}
    null_indices::Vector{Int}
end

struct PilotConfig
    pilot_indices::Vector{Int}
    pilot_symbols::Vector{ComplexF64}
    data_indices::Vector{Int}
end

struct Interleaver
    pi::Vector{Int}      # interleaved position i receives original bit pi[i]
    invpi::Vector{Int}   # original bit j maps to interleaved position invpi[j]
end

struct CarrierBitMap
    data_indices::Vector{Int}
    carrier_to_interleaved_bits::Dict{Int,Tuple{Int,Int}}
    carrier_to_original_bits::Dict{Int,Tuple{Int,Int}}
end

struct LDPCCode
    H::SparseMatrixCSC{Bool,Int}
    A::SparseMatrixCSC{Bool,Int}  # H = [A I]
    n::Int
    k::Int
    m::Int
    var_to_checks::Vector{Vector{Int}}
    check_to_vars::Vector{Vector{Int}}
end

struct BPConfig
    maxiter::Int
    clip::Float64
    algorithm::Symbol # :sumproduct only in v1
end

struct BPResult
    bits::BitVector
    posterior_llr::Vector{Float64}
    success::Bool
    iterations::Int
    syndrome_weight::Int
end
\end{lstlisting}

\begin{lstlisting}
struct RxObservations
    y::Vector{ComplexF64}
    Yfull::Vector{ComplexF64}
    Ypartial_by_M::Dict{Int,Matrix{ComplexF64}}
    noise_power::Float64
end

struct JUNAWeights
    lambda_d::Float64
    eta::Float64
    rho_max::Float64
    gamma_C::Float64
    gamma_W::Float64
    gamma_S::Float64
    lambda_p_max::Float64
    T_rho::Float64
    T_parity::Float64
end

struct JUNAConfig
    M::Int
    Kici::Int
    weights::JUNAWeights
    max_outer::Int
    z_steps_per_outer::Int
    z_stepsize::Float64
    z_clip::Float64
    llr_clip::Float64
    rel_erase_threshold::Float64
    candidate_name::Symbol
end

mutable struct JUNAState
    C::Array{ComplexF64,3}
    W::Matrix{ComplexF64}
    z::Vector{Float64}
    xi::Vector{Float64}
    s::Vector{ComplexF64}
    loss_history::Vector{Float64}
    syndrome_history::Vector{Int}
end
\end{lstlisting}

\clearpage
\section{OFDM, pilots, and partial FFT}

\subsection{OFDM normalization}

The OFDM functions use:
\begin{equation}
x[n]=\frac{1}{\sqrt{N}}\sum_{k=1}^N X[k]e^{j2\pi(k-1)n/N},\quad n=0,\ldots,N-1.
\end{equation}
In Julia:
\begin{lstlisting}
function ofdm_modulate(X::Vector{ComplexF64}, cfg::OFDMConfig)::Vector{ComplexF64}
    @assert length(X) == cfg.N
    return sqrt(cfg.N) .* ifft(X)
end

function ofdm_demodulate(y::Vector{ComplexF64}, cfg::OFDMConfig)::Vector{ComplexF64}
    @assert length(y) == cfg.N
    return fft(y) ./ sqrt(cfg.N)
end
\end{lstlisting}

\subsection{Cyclic prefix}

\begin{lstlisting}
function add_cp(x::Vector{ComplexF64}, cfg::OFDMConfig)::Vector{ComplexF64}
    @assert length(x) == cfg.N
    return vcat(x[end-cfg.Ncp+1:end], x)
end

function remove_cp(ycp::Vector{ComplexF64}, cfg::OFDMConfig)::Vector{ComplexF64}
    @assert length(ycp) == cfg.N + cfg.Ncp
    return ycp[cfg.Ncp+1:end]
end
\end{lstlisting}

\subsection{Pilot selection}

In active-carrier order, every \code{PILOT_SPACING}-th active carrier is a pilot. All other active carriers are data carriers. Null carriers remain zero and are never pilots.

\begin{lstlisting}
function make_comb_pilots(cfg::OFDMConfig, spacing::Int)::PilotConfig
    active = sort(cfg.active_indices)
    pilot_indices = active[1:spacing:end]
    pilot_set = Set(pilot_indices)
    data_indices = [k for k in active if !(k in pilot_set)]
    pilot_symbols = ones(ComplexF64, length(pilot_indices))
    return PilotConfig(pilot_indices, pilot_symbols, data_indices)
end
\end{lstlisting}

\subsection{Deterministic pilot holdout}

If $|\mathcal P|\ge 8\cdot \code{MIN_HOLDOUT_PILOTS}$, validation pilots are every eighth entry in the pilot list. Otherwise, no pilot is held out; validation residual is reported as \code{NaN} and scoring ignores it.

\begin{lstlisting}
function split_pilots_holdout(pcfg::PilotConfig; stride::Int=8, min_holdout::Int=4)
    P = length(pcfg.pilot_indices)
    if P < stride * min_holdout
        return (fit_idx = pcfg.pilot_indices,
                fit_sym = pcfg.pilot_symbols,
                val_idx = Int[],
                val_sym = ComplexF64[])
    end
    val_positions = Set(1:stride:P)
    fit_idx = Int[]; fit_sym = ComplexF64[]
    val_idx = Int[]; val_sym = ComplexF64[]
    for j in 1:P
        if j in val_positions
            push!(val_idx, pcfg.pilot_indices[j]); push!(val_sym, pcfg.pilot_symbols[j])
        else
            push!(fit_idx, pcfg.pilot_indices[j]); push!(fit_sym, pcfg.pilot_symbols[j])
        end
    end
    return (fit_idx=fit_idx, fit_sym=fit_sym, val_idx=val_idx, val_sym=val_sym)
end
\end{lstlisting}

\subsection{Partial FFT}

For $M$ branches, branch $m$ owns samples
\begin{equation}
\mathcal T_m=\left\{n:\left\lfloor\frac{m-1}{M}N\right\rfloor\le n\le \left\lfloor\frac{m}{M}N\right\rfloor-1\right\},
\end{equation}
using 0-based sample index $n$. The rectangular window is
\begin{equation}
g_m[n]=\begin{cases}1,& n\in\mathcal T_m,\\0,&\text{otherwise.}\end{cases}
\end{equation}
Thus $\sum_m g_m[n]=1$.

\begin{equation}
Y_{k,m}=\frac{1}{\sqrt N}\sum_{n=0}^{N-1}g_m[n]y[n]e^{-j2\pi(k-1)n/N}.
\end{equation}

\begin{lstlisting}
function rectangular_partial_windows(N::Int, M::Int)::Matrix{Float64}
    @assert M >= 1
    G = zeros(Float64, M, N)
    for m in 1:M
        n0 = fld((m-1)*N, M) + 1  # Julia 1-based start
        n1 = fld(m*N, M)          # Julia 1-based inclusive end
        G[m, n0:n1] .= 1.0
    end
    @assert maximum(abs.(sum(G, dims=1) .- 1.0)) == 0.0
    return G
end

function partial_fft_observations(y::Vector{ComplexF64}, cfg::OFDMConfig, M::Int)::Matrix{ComplexF64}
    N = cfg.N
    @assert length(y) == N
    G = rectangular_partial_windows(N, M)
    Y = zeros(ComplexF64, N, M)
    for m in 1:M
        Y[:, m] .= fft(G[m, :] .* y) ./ sqrt(N)
    end
    return Y
end
\end{lstlisting}

\textbf{Partial FFT identity test:}
\begin{equation}
\sum_{m=1}^M Y_{k,m}=Y_k^{\rm full}
\end{equation}
for every $k$ and every deterministic test vector $y$.

\clearpage
\section{Interleaver and carrier-to-bit mapping}

The LDPC codeword length must match the number of QPSK data bits:
\begin{equation}
n_b=2|\mathcal D|.
\end{equation}
If this condition fails, \code{make_carrier_bit_map} must throw an error. No silent padding or truncation is allowed in v1.

\subsection{Stride interleaver}

For a codeword of length $n_b$, choose stride $s$ with $\gcd(s,n_b)=1$. The default is $s=33$, but if $\gcd(33,n_b)\ne1$, the constructor searches upward for the first odd stride $s$ with $\gcd(s,n_b)=1$ and logs the choice.

\begin{equation}
\pi[i]=1+((i-1)s\bmod n_b),\qquad i=1,\ldots,n_b.
\end{equation}
The interleaved codeword is
\begin{equation}
b_{\rm int}[i]=b[\pi[i]].
\end{equation}
The inverse satisfies
\begin{equation}
b[j]=b_{\rm int}[\pi^{-1}[j]].
\end{equation}

\begin{lstlisting}
function make_stride_interleaver(n::Int, stride::Int=33)::Interleaver
    s = stride
    while gcd(s, n) != 1
        s += 2
    end
    pi = [1 + mod((i-1)*s, n) for i in 1:n]
    invpi = zeros(Int, n)
    for i in 1:n
        invpi[pi[i]] = i
    end
    return Interleaver(pi, invpi)
end

function interleave_bits(b::BitVector, I::Interleaver)::BitVector
    @assert length(b) == length(I.pi)
    return BitVector([b[I.pi[i]] for i in 1:length(b)])
end

function deinterleave_llr(llr_int::Vector{Float64}, I::Interleaver)::Vector{Float64}
    n = length(llr_int)
    llr = zeros(Float64, n)
    for i in 1:n
        llr[I.pi[i]] = llr_int[i]
    end
    return llr
end
\end{lstlisting}

\subsection{Carrier-bit map}

Sorted data carrier list $\mathcal D=[d_1,\ldots,d_{|\mathcal D|}]$. Carrier $d_r$ carries interleaved bit positions $(2r-1,2r)$. The original LDPC bit positions are $(\pi[2r-1],\pi[2r])$.

\begin{lstlisting}
function make_carrier_bit_map(data_indices::Vector{Int}, code::LDPCCode, I::Interleaver)::CarrierBitMap
    @assert 2 * length(data_indices) == code.n
    @assert length(I.pi) == code.n
    c2i = Dict{Int,Tuple{Int,Int}}()
    c2o = Dict{Int,Tuple{Int,Int}}()
    for (r, k) in enumerate(sort(data_indices))
        i1 = 2r - 1
        i2 = 2r
        c2i[k] = (i1, i2)
        c2o[k] = (I.pi[i1], I.pi[i2])
    end
    return CarrierBitMap(sort(data_indices), c2i, c2o)
end
\end{lstlisting}

\clearpage
\section{LDPC file schema and BP decoder}

\subsection{Systematic TOML schema}

Version 1 accepts only parity-check matrices of the form
\begin{equation}
H=[A\ I_m],\qquad H\in\F_2^{m\times n},\quad n=k+m.
\end{equation}
The codeword is
\begin{equation}
b=[u\ p],\qquad p=A u\pmod 2.
\end{equation}
Then
\begin{equation}
Hb^T=A u+p=0\pmod2.
\end{equation}

Required TOML file format:
\begin{lstlisting}[language={}]
format = "systematic_sparse_v1"
n = 2048
k = 1024
m = 1024
# A is m-by-k. Entries are 1-based row/column coordinates where A[row,col]=1.
A_row = [1, 1, 2, 4]
A_col = [3, 99, 5, 7]
\end{lstlisting}

\begin{lstlisting}
function load_systematic_ldpc_toml(path::String)::LDPCCode
    # parse TOML; verify format, n == k + m; build sparse Bool A
    # H = [A I]
    # build var_to_checks and check_to_vars
end

function ldpc_encode(u::BitVector, code::LDPCCode)::BitVector
    @assert length(u) == code.k
    p = falses(code.m)
    for col in 1:code.k
        if u[col]
            rows = findnz(code.A[:, col])[1]
            for r in rows
                p[r] = xor(p[r], true)
            end
        end
    end
    return vcat(u, p)
end
\end{lstlisting}

\subsection{BP decoder}

Input LLR convention:
\begin{equation}
L_i=\log\frac{P(b_i=0\mid y)}{P(b_i=1\mid y)}.
\end{equation}
Hard decision:
\begin{equation}
\hat b_i=\begin{cases}0,& L_i\ge0,\\1,&L_i<0.\end{cases}
\end{equation}
Check update:
\begin{equation}
m_{j\to i}=2\tanh^{-1}\left(\prod_{i'\in\mathcal N(j)\setminus i}\tanh\frac{m_{i'\to j}}{2}\right),
\end{equation}
with clipping before and after each update. Variable update:
\begin{equation}
m_{i\to j}=L_i+\sum_{j'\in\mathcal N(i)\setminus j}m_{j'\to i}.
\end{equation}
Posterior:
\begin{equation}
L_i^{\rm post}=L_i+\sum_{j\in\mathcal N(i)}m_{j\to i}.
\end{equation}

\begin{lstlisting}
function bp_decode(llr::Vector{Float64}, code::LDPCCode, cfg::BPConfig)::BPResult
    @assert length(llr) == code.n
    # implement flooding schedule sum-product only
    # clip channel input and all messages to cfg.clip
    # stop early if syndrome weight is zero
end
\end{lstlisting}

\clearpage
\section{Channel simulation}

\subsection{Types}

\begin{lstlisting}
struct UWATap
    delay_s::Float64
    power::Float64
    phase::Float64
end

struct DopplerProfile
    a0::Float64
    per_tap_offsets::Vector{Float64}
    slopes::Vector{Float64}
    jitter_std::Float64
end

struct TimeVaryingChannelConfig
    fs::Float64
    taps::Vector{UWATap}
    doppler::DopplerProfile
    rho_h::Float64
    fractional_delay::Symbol # :linear only in v1
end

struct ChannelTruth
    a_true::Float64
    Hfd::Union{Nothing,Matrix{ComplexF64}}
    path_dopplers::Union{Nothing,Matrix{Float64}}
end
\end{lstlisting}

\subsection{Discrete channel model}

For received sample index $n=0,\ldots,N_s-1$,
\begin{equation}
r[n]=\sum_{\ell=1}^{L}\beta_\ell[n]x(n-d_\ell[n])+\nu[n].
\end{equation}
The fractional delay is linear interpolation in v1:
\begin{equation}
x(n-d)=(1-\alpha)x[n_0]+\alpha x[n_0+1],\quad n_0=\lfloor n-d\rfloor,
\quad \alpha=n-d-n_0,
\end{equation}
returning zero for indices outside the array.

Delay evolves by
\begin{equation}
d_\ell[n+1]=d_\ell[n]-a_\ell[n].
\end{equation}
Doppler profile:
\begin{equation}
a_\ell[n]=a_0+\Delta a_\ell+\dot a_\ell\frac{n}{N_s-1}+\epsilon_{\ell,n}.
\end{equation}
For the deterministic global-Doppler tests, use $\Delta a_\ell=\dot a_\ell=\epsilon_{\ell,n}=0$.

\begin{lstlisting}
function apply_time_varying_channel(
    x::Vector{ComplexF64},
    ch::TimeVaryingChannelConfig,
    rng::AbstractRNG;
    add_noise::Bool=false,
    snr_db::Float64=Inf
)::Tuple{Vector{ComplexF64},ChannelTruth}
    # v1 linear fractional delay; optionally add AWGN outside this function in Simulation.jl
end
\end{lstlisting}

\subsection{Channel tests}

\begin{enumerate}
\item \textbf{Delta test.} Input $x[1]=1$, others zero; one tap with integer delay $d=3$, unit gain. Output must be one at sample $4$ and zero elsewhere.
\item \textbf{Fractional delay test.} Input ramp $x[n]=n$ and delay $d=2.5$. Verify linear interpolation against values computed by hand.
\item \textbf{No-Doppler OFDM diagonal test.} A static CP-covered channel must produce negligible off-diagonal empirical $H_{kq}$ when the test applies one basis vector at a time.
\end{enumerate}

\clearpage
\section{Known-Doppler compensation benchmark}

This benchmark reads the truth supplied by the simulation. SC-JUNA never receives \code{ChannelTruth}.

\subsection{Global time scaling}

If the simulator applies global baseband time scaling
\begin{equation}
y[n]=x((1+a)n),
\end{equation}
then compensation samples
\begin{equation}
\tilde y[n]=y\left(\frac{n}{1+a}\right).
\end{equation}
If passband carrier phase is included, multiply by
\begin{equation}
\exp\left[-j2\pi f_c\frac{n}{F_s}\left(1-\frac{1}{1+a}\right)\right].
\end{equation}
For baseband tests in this version, set \code{use_carrier_phase=false}.

\begin{lstlisting}
function compensate_known_doppler(
    y::Vector{ComplexF64},
    a_true::Float64,
    fs::Float64,
    fc::Float64;
    use_carrier_phase::Bool=false
)::Vector{ComplexF64}
    # ytilde[n] = interp_linear(y, n/(1+a_true)) using zero-based n
end
\end{lstlisting}

\subsection{Known-Doppler partial FFT branch phase}

After CP removal, branch midpoint time uses zero-based sample indices:
\begin{equation}
t_m=\frac{1}{F_s}\cdot \frac{n_{m,0}+n_{m,1}}{2},
\end{equation}
where $n_{m,0}=\lfloor(m-1)N/M\rfloor$ and $n_{m,1}=\lfloor mN/M\rfloor-1$.

For subcarrier $k$,
\begin{equation}
d_{k,m}(a)=e^{-j2\pi a f_k t_m}.
\end{equation}
The known-Doppler partial-FFT combined statistic is
\begin{equation}
\chi_k(a)=\frac{1}{M}\sum_{m=1}^M d_{k,m}(a)Y_{k,m}.
\end{equation}

\begin{lstlisting}
function branch_midpoint_time(m::Int, N::Int, M::Int, fs::Float64)::Float64
    n0 = fld((m-1)*N, M)
    n1 = fld(m*N, M) - 1
    return ((n0 + n1) / 2) / fs
end

function known_doppler_partial_fft_combine(
    Ypartial::Matrix{ComplexF64},
    a_true::Float64,
    cfg::OFDMConfig
)::Vector{ComplexF64}
    N, M = size(Ypartial)
    @assert N == cfg.N
    chi = zeros(ComplexF64, N)
    for k in 1:N
        fk = physical_subcarrier_frequency(k, cfg.N, cfg.fs, cfg.fc)
        for m in 1:M
            tm = branch_midpoint_time(m, cfg.N, M, cfg.fs)
            chi[k] += exp(-2im*pi*a_true*fk*tm) * Ypartial[k,m]
        end
        chi[k] /= M
    end
    return chi
end
\end{lstlisting}

\subsection{Known-Doppler undo test}

Use a packet $x$ without a channel or noise. Apply a known synthetic global time scale $a=5\times10^{-5}$ with the same interpolation convention. Then run \code{compensate_known_doppler}. Verify
\begin{equation}
\frac{\norm{\tilde x-x}_2}{\norm{x}_2}<10^{-3}
\end{equation}
for an oversampled deterministic chirp and
\begin{equation}
\frac{\norm{\tilde Y-Y}_2}{\norm{Y}_2}<10^{-2}
\end{equation}
for an OFDM grid. The OFDM tolerance is looser because interpolation is not ideal.

\clearpage
\section{Receiver API}

Each receiver implements this common API.

\begin{lstlisting}
abstract type AbstractReceiver end

struct DecodeResult
    bits::BitVector
    llr::Vector{Float64}
    bp::BPResult
    method::String
    diagnostics::Dict{String,Any}
end

function decode_packet(
    rx::AbstractReceiver,
    ycp::Vector{ComplexF64},
    problem::ReceiverProblem,
    rng::AbstractRNG
)::DecodeResult
    error("receiver not implemented")
end
\end{lstlisting}

\subsection{ReceiverProblem}

\begin{lstlisting}
struct ReceiverProblem
    ofdm::OFDMConfig
    pilots_fit_idx::Vector{Int}
    pilots_fit_sym::Vector{ComplexF64}
    pilots_val_idx::Vector{Int}
    pilots_val_sym::Vector{ComplexF64}
    data_indices::Vector{Int}
    null_indices::Vector{Int}
    interleaver::Interleaver
    bitmap::CarrierBitMap
    ldpc::LDPCCode
    bp::BPConfig
    truth::Union{Nothing,ChannelTruth} # only oracle receivers use this
end
\end{lstlisting}

\subsection{Full-FFT receiver}

\begin{enumerate}
\item Remove CP.
\item Compute \code{Yfull=fft(y)/sqrt(N)}.
\item Estimate $H_p=Y_p/P_p$ on fit pilots.
\item Interpolate $H_k$ by nearest-neighbor in v1.
\item One-tap LMMSE equalize.
\item Estimate effective noise from pilot equalization residuals.
\item Form data LLRs, deinterleave, BP decode.
\end{enumerate}

\begin{equation}
\hat X_k=\frac{\hat H_k^*}{|\hat H_k|^2+\hat\sigma^2/E_s}Y_k.
\end{equation}

\begin{lstlisting}
struct FullFFTReceiver <: AbstractReceiver
    bp::BPConfig
end
\end{lstlisting}

\subsection{RPC receiver}

For target carrier $k$, choose fit pilots within radius \code{local_pilot_radius}; if none exist, use the nearest two pilots. The combiner solves
\begin{equation}
\mathbf w_k=(U_k^HU_k+\delta I)^{-1}U_k^H\mathbf d_k.
\end{equation}

\begin{lstlisting}
struct RPCConfig
    M::Int
    delta::Float64
    local_pilot_radius::Int
end

struct RPCReceiver <: AbstractReceiver
    cfg::RPCConfig
    bp::BPConfig
end

function fit_rpc_weights(
    Ypartial::Matrix{ComplexF64},
    pilot_idx::Vector{Int},
    pilot_sym::Vector{ComplexF64},
    cfg::RPCConfig
)::Matrix{ComplexF64}
end
\end{lstlisting}

\subsection{Known-Doppler OFDM receiver}

Requires \code{problem.truth !== nothing}. It compensates with true global Doppler, then runs the full-FFT receiver. It must throw an error if truth is missing.

\begin{lstlisting}
struct KnownDopplerOFDMReceiver <: AbstractReceiver
    use_carrier_phase::Bool
    bp::BPConfig
end
\end{lstlisting}

\subsection{Known-Doppler partial-FFT receiver}

Requires \code{problem.truth !== nothing}. It computes partial FFT, combines branches with true Doppler phase, then runs pilot channel estimation and LDPC.

\begin{lstlisting}
struct KnownDopplerPartialFFTReceiver <: AbstractReceiver
    M::Int
    bp::BPConfig
end
\end{lstlisting}

\clearpage
\section{Packet diagnostics and candidate list}

\subsection{Packet diagnostics}

Diagnostics use only the current packet. No feedback and no truth object.

\begin{align}
\widehat\sigma^2&=\operatorname{mean}_{k\in\mathcal N}|Y_k|^2\quad\text{if nulls exist, else pilot residual median},\\
\widehat H_p&=Y_p/P_p,\\
\operatorname{SFM}&=\frac{\exp\left(\frac{1}{|\mathcal P|}\sum_p\log(|\hat H_p|^2+\eps)\right)}{\frac{1}{|\mathcal P|}\sum_p(|\hat H_p|^2+\eps)},\\
D_{\rm notch}&=10\log_{10}\frac{\max_p |\hat H_p|^2+\eps}{\min_p |\hat H_p|^2+\eps},\\
\eta_{\rm br}&=\frac{\sum_{p,m}|Y_{p,m}-\bar Y_p|^2}{\sum_{p,m}|Y_{p,m}|^2+\eps}.
\end{align}

\begin{lstlisting}
struct PacketDiagnostics
    snr_db::Float64
    noise_power::Float64
    sfm::Float64
    notch_depth_db::Float64
    branch_disagreement::Float64
    rpc_pilot_residual::Float64
end
\end{lstlisting}

\subsection{Base $M_0$ and $K_0$ rules}

Let $\eta_b=\eta_{\rm br}$ and $s=\operatorname{SFM}$.
\begin{equation}
M_0=\begin{cases}
4,&\eta_b<0.08,\\
8,&0.08\le\eta_b<0.20,\\
16,&0.20\le\eta_b<0.40,\\
32,&\eta_b\ge0.40.
\end{cases}
\end{equation}

\begin{equation}
K_0=\begin{cases}
0,&\eta_b<0.05\text{ and }s>0.60,\\
1,&\eta_b<0.15,\\
2,&0.15\le\eta_b<0.35,\\
3,&\eta_b\ge0.35.
\end{cases}
\end{equation}
If $D_{\rm notch}>25$ dB, set $K_0\leftarrow\min(K_0+1,K_{\max})$.

\subsection{Candidate list}

Always generate candidates in this order:
\begin{equation}
[\code{conservative},\code{nominal},\code{doppler_heavy},\code{notchy},\code{low_snr},\code{rescue}].
\end{equation}
Candidate modifications are fixed in Section~\ref{sec:constants}. Do not invent extra candidates.

\begin{lstlisting}
function build_scjuna_candidates(diag::PacketDiagnostics, base::JUNAConfig)::Vector{JUNAConfig}
    M0 = base_M_rule(diag.branch_disagreement)
    K0 = base_K_rule(diag.branch_disagreement, diag.sfm, diag.notch_depth_db)
    return [make_candidate(:conservative, base, M0, K0),
            make_candidate(:nominal, base, M0, K0),
            make_candidate(:doppler_heavy, base, M0, K0),
            make_candidate(:notchy, base, M0, K0),
            make_candidate(:low_snr, base, M0, K0),
            make_candidate(:rescue, base, M0, K0)]
end
\end{lstlisting}

\section{Initialization}

\begin{enumerate}
\item Fit RPC weights using fit pilots only.
\item Initialize $W^{(0)}$ from RPC.
\item Initialize $C^{(0)}$ with only diagonal local branch coefficient nonzero:
\begin{equation}
c_{kk,m}^{(0)}=1/M,\\qquad c_{kq,m}^{(0)}=0 \text{ for } q\ne k.
\end{equation}
\item Compute $\tilde s_k=(w_k^{(0)})^H Y_k$.
\item Convert $\tilde s_k$ to initial LLRs; deinterleave; initialize
\begin{equation}
z_i^{(0)}=\operatorname{clip}(L_i/2,-2,2).
\end{equation}
\item Compute $\xi_i=\tanh z_i$ and relaxed symbols.
\end{enumerate}

\begin{lstlisting}
function initialize_juna_state(
    Ypartial::Matrix{ComplexF64},
    problem::ReceiverProblem,
    cfg::JUNAConfig,
    rpc::RPCConfig
)::JUNAState
    # exact steps above
end
\end{lstlisting}

\clearpage
\section{JUNA objective}

For candidate configuration $\theta=(M,K,\lambda_d,\eta,\rho,\gamma_C,\gamma_W,\gamma_S,\lambda_p)$, define branch residual
\begin{equation}
\mathbf r_k=\mathbf Y_k-\sum_{q\in\mathcal B_K(k)}\mathbf c_{kq}s_q.
\end{equation}
Scalar combiner output:
\begin{equation}
\hat s_k=\mathbf w_k^H\mathbf Y_k.
\end{equation}
Relaxed bits:
\begin{equation}
\xi_i=\tanh z_i.
\end{equation}
For data carrier $k$ with original LDPC bit indices $(i_I,i_Q)$,
\begin{equation}
s_k(z)=\frac{1}{\sqrt2}(\xi_{i_I}+j\xi_{i_Q}).
\end{equation}
For pilots, $s_p=P_p$; for nulls, $s_k=0$.

The loss is
\begin{align}
\mathcal L &=
\lambda_d\sum_{k=1}^N\norm{\mathbf r_k}_2^2
+\eta\sum_{p\in\mathcal P_{\rm fit}}|\mathbf w_p^H\mathbf Y_p-P_p|^2
+\rho^{(t)}\sum_{k\in\mathcal D}|\mathbf w_k^H\mathbf Y_k-s_k|^2\nonumber\\
&\quad+\gamma_C\norm{C}_F^2
+\gamma_W\sum_{k=1}^N\norm{\mathbf w_k}_2^2
+\gamma_S\sum_{k=1}^{N-1}\norm{\mathbf w_{k+1}-\mathbf w_k}_2^2
+\lambda_p^{(t)}\mathcal R_{\rm par}(z),
\end{align}
where
\begin{equation}
\mathcal R_{\rm par}(z)=\sum_{j=1}^{m}\left(1-\prod_{i\in\mathcal N_j}\xi_i\right)^2.
\end{equation}

\clearpage
\section{JUNA block updates}

\subsection{$C$ update}

For each received carrier $k$, let $\mathbf s_k=[s_q]_{q\in\mathcal B_K(k)}$. Let $C_k\in\C^{M\times B_k}$ collect valid branch coefficients. Solve
\begin{equation}
C_k=\arg\min_C\lambda_d\norm{\mathbf Y_k-C\mathbf s_k}_2^2+\gamma_C\norm C_F^2.
\end{equation}
Closed form:
\begin{equation}
C_k=\lambda_d\mathbf Y_k\mathbf s_k^H\left(\lambda_d\mathbf s_k\mathbf s_k^H+\gamma_C I\right)^{-1}.
\end{equation}
Write valid coefficients into \code{C[k, off, m]}; leave invalid edge slots zero.

\subsection{$W$ update}

Without smoothness, use
\begin{equation}
\mathbf w_k=\left(a_k\mathbf Y_k\mathbf Y_k^H+\gamma_W I\right)^{-1}a_k\mathbf Y_k d_k^*,
\end{equation}
where
\begin{equation}
(a_k,d_k)=\begin{cases}
(\eta,P_k),&k\in\mathcal P_{\rm fit},\\
(\rho^{(t)},s_k),&k\in\mathcal D,\\
(0,0),&\text{otherwise.}
\end{cases}
\end{equation}
For v1, smoothness $\gamma_S$ is applied by one smoothing pass after each update:
\begin{equation}
\mathbf w_k\leftarrow \frac{\mathbf w_k+\alpha_s(\mathbf w_{k-1}+\mathbf w_{k+1})}{1+2\alpha_s},
\qquad
\alpha_s=\min(0.25,\gamma_S).
\end{equation}
At edges, use the available neighbor only. This avoids a block-tridiagonal implementation in v1. The fixed smoothing pass after each update applies $\gamma_S$. A later version may replace this with the exact block-tridiagonal solve.

\subsection{$z$ update}

For each outer iteration, perform \code{z_steps_per_outer} gradient steps:
\begin{equation}
z\leftarrow \operatorname{clip}(z-\mu g_z,-z_{\max},z_{\max}).
\end{equation}
Then recompute $\xi=\tanh z$ and all data symbols $s_k$.

\begin{lstlisting}
function run_juna_solver(
    init::JUNAState,
    obs::RxObservations,
    problem::ReceiverProblem,
    cfg::JUNAConfig
)::JUNAResult
    state = init
    Yp = obs.Ypartial_by_M[cfg.M]
    grad_z = zeros(Float64, length(state.z))
    for t in 0:cfg.max_outer-1
        rho_t = cfg.weights.rho_max * (1 - exp(-(t+1)/cfg.weights.T_rho))
        lambda_p_t = cfg.weights.lambda_p_max * (1 - exp(-(t+1)/cfg.weights.T_parity))
        update_C!(state, Yp, problem, cfg)
        update_W!(state, Yp, problem, cfg, rho_t)
        for step in 1:cfg.z_steps_per_outer
            juna_z_gradient!(grad_z, state, Yp, problem, cfg, rho_t, lambda_p_t)
            state.z .-= cfg.z_stepsize .* grad_z
            clamp!(state.z, -cfg.z_clip, cfg.z_clip)
            update_xi_and_symbols!(state, problem)
        end
        llr = juna_llrs(state, Yp, problem, cfg)
        bp = bp_decode(llr, problem.ldpc, problem.bp)
        push!(state.syndrome_history, bp.syndrome_weight)
        push!(state.loss_history, juna_loss(state, Yp, problem, cfg, rho_t, lambda_p_t))
        if bp.success
            return JUNAResult(bp.bits, llr, bp, state, cfg, nothing, Dict{String,Any}())
        end
    end
    llr = juna_llrs(state, Yp, problem, cfg)
    bp = bp_decode(llr, problem.ldpc, problem.bp)
    return JUNAResult(bp.bits, llr, bp, state, cfg, nothing, Dict{String,Any}())
end
\end{lstlisting}

\section{Manual gradient for the $z$ update}

Production code uses the hand gradient below. AD is allowed only in tests to check this gradient.

\subsection{Symbol-gradient accumulation}

Define complex symbol gradient
\begin{equation}
g_s[q]=\frac{\partial\mathcal L}{\partial \Re s_q}+j\frac{\partial\mathcal L}{\partial \Im s_q}.
\end{equation}
Initialize $g_s[q]=0$ for all data carriers.

Observation term: for every received carrier $k$ and valid neighbor $q\in\mathcal B_K(k)$,
\begin{equation}
u_{kq}=\mathbf c_{kq}^H\mathbf r_k,
\qquad
 g_s[q]\leftarrow g_s[q]-2\lambda_d u_{kq}.
\end{equation}
Combiner-symbol tie for data carrier $q$:
\begin{equation}
e_q=\mathbf w_q^H\mathbf Y_q-s_q,
\qquad
 g_s[q]\leftarrow g_s[q]-2\rho^{(t)}e_q.
\end{equation}

\subsection{Convert symbol gradient to $z$ gradient}

For carrier $q$ with original bit indices $(i_I,i_Q)$,
\begin{align}
g_\xi[i_I]&\leftarrow g_\xi[i_I]+\frac{\Re g_s[q]}{\sqrt2},\\
g_\xi[i_Q]&\leftarrow g_\xi[i_Q]+\frac{\Im g_s[q]}{\sqrt2}.
\end{align}

Parity term: for each check $j$, define
\begin{equation}
P_j=\prod_{i\in\mathcal N_j}\xi_i.
\end{equation}
For each $i\in\mathcal N_j$,
\begin{equation}
g_\xi[i]\leftarrow g_\xi[i]-2\lambda_p^{(t)}(1-P_j)
\prod_{r\in\mathcal N_j\setminus i}\xi_r.
\end{equation}
Do not use division by $\xi_i$; use the product excluding $i$ to avoid zero problems.

Finally,
\begin{equation}
\frac{\partial\mathcal L}{\partial z_i}=g_\xi[i](1-\xi_i^2).
\end{equation}

\begin{lstlisting}
function juna_z_gradient!(
    grad_z::Vector{Float64},
    state::JUNAState,
    Ypartial::Matrix{ComplexF64},
    problem::ReceiverProblem,
    cfg::JUNAConfig,
    rho_t::Float64,
    lambda_p_t::Float64
)::Nothing
    # 1. zero grad_z, g_xi, g_s
    # 2. compute all residual r_k
    # 3. accumulate observation gradient into g_s
    # 4. accumulate combiner-symbol tie into g_s
    # 5. map g_s to g_xi through CarrierBitMap
    # 6. add parity gradient using product-excluding-i
    # 7. grad_z[i] = g_xi[i] * (1 - xi[i]^2)
end
\end{lstlisting}

\subsection{Gradient test}

For a tiny deterministic problem $N=16$, $K=1$, $M=4$, and a small LDPC toy code, compare manual gradient to central finite difference:
\begin{equation}
\left|g_i^{\rm manual}-\frac{L(z+h e_i)-L(z-h e_i)}{2h}\right|<10^{-4}
\end{equation}
for $h=10^{-5}$ and all tested coordinates. Test at least 16 coordinates including bits in multiple parity checks.

\clearpage
\section{LLR formation and reliability erasure}

For scalar output $\hat s_k=\mathbf w_k^H\mathbf Y_k$, estimate pilot residuals
\begin{equation}
e_p=|\mathbf w_p^H\mathbf Y_p-P_p|^2.
\end{equation}
For each data carrier $k$, set $\sigma_{\rm eff,k}^2$ to the median of the three nearest pilot residuals plus \code{EPS}. If fewer than three pilots exist, use the median of all pilot residuals. Define raw reliability
\begin{equation}
\gamma_k=\frac{1}{\sigma_{\rm eff,k}^2+\eps}.
\end{equation}
Define relative reliability
\begin{equation}
\gamma_k^{\rm rel}=\frac{\gamma_k}{\operatorname{median}_{q\in\mathcal D}\gamma_q+\eps}.
\end{equation}
If $\gamma_k^{\rm rel}<\code{REL_ERASE_THRESHOLD}$, set both LLRs for carrier $k$ to zero.

For QPSK carriers that are not erased:
\begin{align}
L_{I,k}&=\operatorname{clip}\left(\frac{2\sqrt2\Re\hat s_k}{\sigma_{\rm eff,k}^2},-L_{\max},L_{\max}\right),\\
L_{Q,k}&=\operatorname{clip}\left(\frac{2\sqrt2\Im\hat s_k}{\sigma_{\rm eff,k}^2},-L_{\max},L_{\max}\right).
\end{align}
These fill \code{llr_int[2r-1]} and \code{llr_int[2r]} for data carrier rank $r$. Deinterleave before BP.

\begin{lstlisting}
function juna_llrs(
    state::JUNAState,
    Ypartial::Matrix{ComplexF64},
    problem::ReceiverProblem,
    cfg::JUNAConfig
)::Vector{Float64}
    # compute shat, sigma_eff, relative reliability
    # fill llr_int in carrier order
    # erase unreliable carriers
    # clip and deinterleave
end
\end{lstlisting}

\section{Candidate score}

If validation pilots exist:
\begin{equation}
e_{\rm val}=\frac{\sum_{p\in\mathcal P_{\rm val}}|\mathbf w_p^H\mathbf Y_p-P_p|^2}{\sum_{p\in\mathcal P_{\rm val}}|P_p|^2+\eps}.
\end{equation}
If validation pilots do not exist, set $e_{\rm val}=0$ and do not use it in the score.

Normalize residuals:
\begin{align}
R_{\rm syn}&=\frac{\text{syndrome weight}}{m},\\
R_{\rm model}&=\frac{\sum_k\norm{\mathbf r_k}_2^2}{\sum_k\norm{\mathbf Y_k}_2^2+\eps},\\
R_C&=\frac{\norm C_F^2}{N(2K+1)M+\eps},\\
R_W&=\frac{\norm W_F^2}{NM+\eps}.
\end{align}
Score:
\begin{equation}
\mathcal V=e_{\rm val}+0.5R_{\rm syn}+0.5R_{\rm model}+10^{-3}R_C+10^{-3}R_W.
\end{equation}

Selection priority:
\begin{enumerate}
\item If CRC is enabled and any candidate passes CRC, choose the passing candidate with lowest $e_{\rm val}$, then lowest candidate index.
\item Else if any candidate has BP syndrome zero, choose the passing candidate with lowest $e_{\rm val}$ if validation exists; otherwise lowest score.
\item Else choose lowest score.
\end{enumerate}

\clearpage
\section{SC-JUNA receiver}

\begin{lstlisting}
function scjuna_decode(
    ycp::Vector{ComplexF64},
    problem::ReceiverProblem,
    scfg::SCJUNAConfig,
    rng::AbstractRNG
)::JUNAResult
    # 1. aligned packet only
    y = remove_cp(ycp, problem.ofdm)

    # 2. full FFT and all candidate partial FFTs; recompute each M
    Yfull = ofdm_demodulate(y, problem.ofdm)
    Ydict = Dict{Int,Matrix{ComplexF64}}()
    for M in scfg.candidate_M
        Ydict[M] = partial_fft_observations(y, problem.ofdm, M)
    end

    # 3. noise estimate from nulls or pilot fallback
    noise_power = estimate_noise_power(Yfull, problem.null_indices,
                                       problem.pilots_fit_idx, problem.pilots_fit_sym)
    obs = RxObservations(y, Yfull, Ydict, noise_power)

    # 4. diagnostics use largest M for branch disagreement
    diag = compute_packet_diagnostics(obs, problem, scfg.rpc)

    # 5. deterministic candidate list
    candidates = build_scjuna_candidates(diag, scfg.base_juna)

    # 6. evaluate all candidates in fixed order
    results = JUNAResult[]
    scores = CandidateScore[]
    for cfg in candidates
        Yp = obs.Ypartial_by_M[cfg.M]
        init = initialize_juna_state(Yp, problem, cfg, scfg.rpc)
        result = run_juna_solver(init, obs, problem, cfg)
        score = score_juna_candidate(result, obs, problem, cfg)
        push!(results, result)
        push!(scores, score)
    end

    # 7. select and attach diagnostics
    ibest = select_best_candidate(results, scores)
    best = results[ibest]
    best.score = scores[ibest]
    best.diagnostics["packet"] = diag
    best.diagnostics["candidate_scores"] = scores
    return best
end
\end{lstlisting}

\clearpage
\section{Tests}

These tests must be run before any Monte Carlo experiment. Each test uses fixed seeds and small dimensions.

\subsection{Indexing tests}

For $N=8$:
\begin{equation}
[\kappa(1),\ldots,\kappa(8)]=[0,1,2,3,4,-3,-2,-1].
\end{equation}
For $k=1$, $K=2$, valid neighbors are $[1,2,3]$ and offsets are $[3,4,5]$. Invalid offsets 1 and 2 are zero slots.

\subsection{OFDM roundtrip}

Use $N=16$, $N_{cp}=4$, deterministic QPSK grid. Verify
\begin{equation}
\norm{\operatorname{FFT}(\operatorname{IFFT}(X))-X}_\infty<10^{-12}.
\end{equation}

\subsection{Partial FFT identity}

Use $N=16$, $M=4$, deterministic complex vector $y[n]=(n+1)+j(2n-1)$. Verify
\begin{equation}
\max_k\left|\sum_mY_{k,m}-Y_k^{\rm full}\right|<10^{-12}.
\end{equation}

\subsection{LDPC encode and parity}

Toy code:
\begin{equation}
A=\begin{bmatrix}1&0&1\\0&1&1\end{bmatrix},\qquad H=[A\ I_2].
\end{equation}
For $u=[1,0,1]$, $p=A u=[0,1]$ over $\F_2$ and $b=[1,0,1,0,1]$. Verify $Hb^T=0$.

\subsection{BP noiseless decode}

For the toy code above, use LLRs $L_i=+20$ for bit 0 and $-20$ for bit 1. BP must succeed in one iteration or less than five iterations.

\subsection{Channel tests}

Run the three tests in the channel section: integer-delay delta, fractional-delay ramp, and static diagonal OFDM channel.

\subsection{Known-Doppler undo test}

Use the required test in the known-Doppler section. This catches sign, unit, and indexing mistakes.

\subsection{RPC synthetic test}

Construct $Y_{k,m}=a_m X_k$ for all carriers with known branch vector $a$. Fit RPC on pilots. Verify $w^Ha\approx1$ and pilot residual below $10^{-10}$ for data without noise.

\subsection{JUNA gradient test}

Use the finite-difference test in the manual gradient section. This test must pass before running SC-JUNA.

\subsection{SC-JUNA candidate selection test}

Create two diagnostics:
\begin{align}
\eta_{\rm br}=0.02,\quad {\rm SFM}=0.9 &\Rightarrow M_0=4, K_0=0,\\
\eta_{\rm br}=0.50,\quad {\rm SFM}=0.4 &\Rightarrow M_0=32, K_0\ge3.
\end{align}
Verify generated candidate list has fixed order and expected $M,K$ trends.

\clearpage
\section{Simulation loop}

\begin{lstlisting}
struct TxPacket
    info_bits::BitVector
    coded_bits::BitVector
    interleaved_bits::BitVector
    X::Vector{ComplexF64}
    x::Vector{ComplexF64}
    xcp::Vector{ComplexF64}
end

function generate_packet(problem::ReceiverProblem, rng::AbstractRNG)::TxPacket
    u = BitVector(rand(rng, Bool, problem.ldpc.k))
    b = ldpc_encode(u, problem.ldpc)
    bint = interleave_bits(b, problem.interleaver)
    sdata = qpsk_mod(bint)
    X = zeros(ComplexF64, problem.ofdm.N)
    # fill data carriers in sorted order using sdata
    # fill fit and validation pilots with known symbols
    x = ofdm_modulate(X, problem.ofdm)
    xcp = add_cp(x, problem.ofdm)
    return TxPacket(u, b, bint, X, x, xcp)
end
\end{lstlisting}

One trial:
\begin{enumerate}
\item Generate packet.
\item Apply channel.
\item Add AWGN with the same noise draw for all receivers at the same SNR.
\item Decode using each receiver.
\item Compare decoded information bits after extracting the first $k$ bits of the systematic LDPC codeword.
\item Log BER, FER, syndrome weight, BP success, selected $M$, selected $K$, diagnostics, and runtime.
\end{enumerate}

\section{Benchmark receiver set}

Minimum benchmark suite:
\begin{enumerate}
\item \code{FullFFTReceiver}: ordinary coherent OFDM with pilot channel estimate.
\item \code{RPCReceiver}: partial-FFT combiner baseline.
\item \code{KnownDopplerOFDMReceiver}: oracle global Doppler compensation plus coherent OFDM.
\item \code{KnownDopplerPartialFFTReceiver}: oracle Doppler partial-FFT baseline.
\item \code{FixedJUNAReceiver}: JUNA with one fixed config.
\item \code{SCJUNAReceiver}: SC-JUNA.
\end{enumerate}
Optional upper bound:
\begin{equation}
\hat X=(H^HH+\sigma^2I/E_s)^{-1}H^HY
\end{equation}
using true $H$ from simulation. Label as genie only.

\section{Done criteria}

The implementation is considered ready for Monte Carlo sweeps only when all conditions hold:
\begin{enumerate}
\item All tests in this specification pass.
\item Full-FFT receiver decodes noiseless static channels with zero bit errors.
\item RPC receiver matches full FFT in a static time-invariant channel when equal gain combining is optimal.
\item Known-Doppler receiver improves or matches uncompensated full FFT in the synthetic global-Doppler test.
\item JUNA gradient test passes against finite differences.
\item SC-JUNA candidate selection is deterministic across repeated runs with the same inputs.
\item No function in the receive path calls global randomness.
\end{enumerate}

\clearpage
\appendix
\section{Background equations}

This appendix is intentionally short. The implementation should rely on the function definitions in the main body, not on deriving these equations again.

For time-varying multipath,
\begin{equation}
r(t)=\sum_{\ell=1}^L\beta_\ell(t)x(t-\tau_\ell(t))+\nu(t).
\end{equation}
After CP removal and FFT,
\begin{equation}
Y_k=\sum_qH_{kq}X_q+\nu_k.
\end{equation}
If the channel is time-invariant over the OFDM symbol, $H_{kq}$ is diagonal. If Doppler or path delays vary within the symbol, off-diagonal ICI appears. Partial FFT preserves observations from each short window
\begin{equation}
Y_{k,m}=\frac{1}{\sqrt N}\sum_{n=0}^{N-1}g_m[n]y[n]e^{-j2\pi(k-1)n/N},
\end{equation}
which JUNA combines with pilots and LDPC constraints.

\section{Research references to consult separately}

The production spec does not require reading these to implement code, but they motivate the design:
\begin{enumerate}
\item S. Yerramalli, M. Stojanovic, and U. Mitra, ``Partial FFT Demodulation: A Detection Method for Highly Doppler Distorted OFDM Systems,'' IEEE Transactions on Signal Processing, 2012.
\item The JUNA project draft defining the coupled branch coefficients, combiner, and LDPC parity objective.
\end{enumerate}

\end{document}
